\begin{center}\begin{tikzpicture}[font=\small,>=stealth,line join=round]
\draw[gray!40,rounded corners=3pt] (0,0) rectangle (14.2,8.3);\node at (7.1,7.700000000000001) {例2：用周期4只能列出偶数轴，会漏掉奇数轴};
\begin{scope}[shift={(2.1714285714285717,3.8222222222222224)},x=1.0714285714285714cm,y=1.277777777777778cm]\draw[->,gray] (-1,0)--(4.6,0) node[right]{$x$};\draw[->,gray] (0,-1.7)--(0,1.9) node[above]{$y$};\draw[blue,thick,solid,domain=-0.7:4.2,samples=180] plot(\x,{cos(deg(pi*\x))});\foreach \a in {0,2,4}{\draw[violet,dashed,thick] (\a,-1.3)--(\a,1.3);}\foreach \a in {1,3}{\draw[orange!85!black,dotted,thick] (\a,-1.3)--(\a,1.3);}\end{scope}\node[anchor=west,align=left,text width=5.0cm] at (8.4,4.0) {$f=\cos\pi x$\\[9pt]已知一个周期 $T=4$\\[9pt]紫线：公式保证 $x=2k$\\[9pt]橙点线：$x=1,3$ 也是真轴\\[9pt]最小周期2才列齐 $x=k$};
\node[text width=12.6cm,align=center] at (7.1,.65) {两类竖线都是真对称轴；已知的一个周期越大，用它列出的一组位置可能越稀。};
\end{tikzpicture}\end{center}
